% ========================================================================
% ILLUSTRATIONS ET CAPTURES MANQUANTES
% ========================================================================

\chapter{Guide des Illustrations et Captures à Ajouter}

\begin{objectivebox}
Ce chapitre liste toutes les illustrations, captures d'écran et diagrammes mentionnés dans le rapport qui doivent être ajoutés pour compléter la documentation.
\end{objectivebox}

\section{Captures d'Écran du Chapitre Introduction}

\subsection{Contexte et Présentation}

\begin{enumerate}
    \item \textbf{Logo et présentation de l'entreprise ILOGsys}
    \begin{itemize}
        \item Localisation : Section 1.1 - Présentation de l'Entreprise
        \item Description : Logo officiel, organigramme, domaines d'expertise
        \item Format suggéré : Image PNG haute résolution
    \end{itemize}
    
    \item \textbf{Diagramme illustrant les problématiques du secteur}
    \begin{itemize}
        \item Localisation : Section 1.1 - Problématique du Secteur
        \item Description : Infographie des défis (fragmentation, communication, estimation, suivi)
        \item Format suggéré : Diagramme vectoriel ou infographie
    \end{itemize}
    
    \item \textbf{Schéma des objectifs et leur interconnexion}
    \begin{itemize}
        \item Localisation : Section 1.2 - Objectifs Spécifiques
        \item Description : Mind map des 5 objectifs principaux avec liaisons
        \item Format suggéré : Diagramme conceptuel
    \end{itemize}
    
    \item \textbf{Diagramme CRISP-DM adapté au projet Housy}
    \begin{itemize}
        \item Localisation : Section 1.3 - Les Six Phases Appliquées
        \item Description : Cycle CRISP-DM avec adaptations pour développement web
        \item Format suggéré : Diagramme circulaire avec flèches
    \end{itemize}
    
    \item \textbf{Tableau récapitulatif du périmètre et contraintes}
    \begin{itemize}
        \item Localisation : Section 1.4 - Contraintes Métier
        \item Description : Tableau structuré des contraintes techniques/métier/fonctionnelles
        \item Format suggéré : Tableau formaté avec codes couleur
    \end{itemize}
\end{enumerate}

\subsection{Processus et Documentation}

\begin{enumerate}
    \item \textbf{Diagramme de classification des documents}
    \begin{itemize}
        \item Localisation : Section "Description des Documents"
        \item Description : Classification hiérarchique des types de documents
        \item Format suggéré : Organigramme ou arbre de classification
    \end{itemize}
    
    \item \textbf{Cartographie de l'écosystème numérique actuel}
    \begin{itemize}
        \item Localisation : Section "Digitalisation de la Construction"
        \item Description : État des lieux des outils numériques utilisés
        \item Format suggéré : Cartographie avec acteurs et outils
    \end{itemize}
    
    \item \textbf{Chronogramme détaillé du processus de construction}
    \begin{itemize}
        \item Localisation : Section "Description Textuelle du Processus"
        \item Description : Timeline des 4 phases avec durées et jalons
        \item Format suggéré : Diagramme de Gantt ou chronogramme
    \end{itemize}
    
    \item \textbf{Diagramme des flux multicouches}
    \begin{itemize}
        \item Localisation : Section "Description des Flux Échangés"
        \item Description : Visualisation des flux d'information, financiers, matériels
        \item Format suggéré : Diagramme en couches avec codes couleur
    \end{itemize}
\end{enumerate}

\subsection{Modélisation et Architecture}

\begin{enumerate}
    \item \textbf{Diagramme BPMN complet du processus de construction}
    \begin{itemize}
        \item Localisation : Section "Modélisation du Processus"
        \item Description : Processus complet en notation BPMN 2.0
        \item Format suggéré : Diagramme BPMN professionnel avec pools et lanes
        \item Statut : ✅ CRÉÉ dans processus\_construction.tex
    \end{itemize}
    
    \item \textbf{Diagramme de cas d'usage UML}
    \begin{itemize}
        \item Localisation : Section "Diagrammes UML Complémentaires"
        \item Description : Cas d'usage principaux avec acteurs
        \item Format suggéré : Diagramme UML standard
        \item Statut : ✅ CRÉÉ dans diagrammes\_uml.tex (existant)
    \end{itemize}
    
    \item \textbf{Diagramme de séquence pour l'estimation de coûts}
    \begin{itemize}
        \item Localisation : Section "Diagrammes UML Complémentaires"
        \item Description : Interactions temporelles estimation IA
        \item Format suggéré : Diagramme de séquence UML
        \item Statut : ✅ CRÉÉ dans processus\_construction.tex
    \end{itemize}
    
    \item \textbf{Diagramme d'activité UML pour le processus d'approbation}
    \begin{itemize}
        \item Localisation : Section "Diagrammes UML Complémentaires"
        \item Description : Workflow validation avec décisions et parallélisme
        \item Format suggéré : Diagramme d'activité UML
        \item Statut : ✅ CRÉÉ dans processus\_construction.tex
    \end{itemize}
\end{enumerate}

\subsection{Solution et Architecture Technique}

\begin{enumerate}
    \item \textbf{Analyse comparative des solutions existantes}
    \begin{itemize}
        \item Localisation : Section "Critiques de l'Existant"
        \item Description : Tableau comparatif ERP/Spécialisés/Maison avec forces/faiblesses
        \item Format suggéré : Tableau matriciel avec scoring
    \end{itemize}
    
    \item \textbf{Architecture technique de Housy Tunisia}
    \begin{itemize}
        \item Localisation : Section "Architecture de la Solution"
        \item Description : Architecture multicouche avec composants techniques
        \item Format suggéré : Diagramme d'architecture système
        \item Statut : ✅ CRÉÉ dans processus\_construction.tex
    \end{itemize}
    
    \item \textbf{Diagramme des modules fonctionnels}
    \begin{itemize}
        \item Localisation : Section "Modules Fonctionnels"
        \item Description : 5 modules principaux avec interconnexions
        \item Format suggéré : Diagramme modulaire avec flux de données
    \end{itemize}
    
    \item \textbf{Roadmap des chapitres suivants}
    \begin{itemize}
        \item Localisation : Section "Conclusion"
        \item Description : Plan visuel des chapitres avec liens CRISP-DM
        \item Format suggéré : Roadmap visuelle ou timeline
    \end{itemize}
\end{enumerate}

\section{Nouvelles Illustrations Créées}

\subsection{Diagrammes Techniques Générés}

\begin{enumerate}
    \item \textbf{Diagramme BPMN Global du Processus de Construction}
    \begin{itemize}
        \item Fichier : \texttt{processus\_construction.tex}
        \item Description : Processus complet avec 5 pools d'acteurs
        \item Éléments : Événements, activités, passerelles, messages
        \item Statut : ✅ COMPLÉTÉ
    \end{itemize}
    
    \item \textbf{Architecture des Flux Multicouches}
    \begin{itemize}
        \item Fichier : \texttt{processus\_construction.tex}
        \item Description : 4 couches de flux (information, financier, matériel, humain)
        \item Éléments : Acteurs connectés, flux colorés par type
        \item Statut : ✅ COMPLÉTÉ
    \end{itemize}
    
    \item \textbf{Diagramme de Séquence - Estimation IA}
    \begin{itemize}
        \item Fichier : \texttt{processus\_construction.tex}
        \item Description : 13 étapes d'interaction pour estimation
        \item Éléments : 6 acteurs, messages temporels, zones d'activation
        \item Statut : ✅ COMPLÉTÉ
    \end{itemize}
    
    \item \textbf{Diagramme d'Activité - Validation}
    \begin{itemize}
        \item Fichier : \texttt{processus\_construction.tex}
        \item Description : Workflow validation avec activités parallèles
        \item Éléments : Fork/join, décisions, zones de responsabilité
        \item Statut : ✅ COMPLÉTÉ
    \end{itemize}
    
    \item \textbf{Architecture Technique Multicouche}
    \begin{itemize}
        \item Fichier : \texttt{processus\_construction.tex}
        \item Description : 6 couches avec composants détaillés
        \item Éléments : Présentation, Application, Métier, IA, Données, Infrastructure
        \item Statut : ✅ COMPLÉTÉ
    \end{itemize}
\end{enumerate}

\section{Captures d'Écran de l'Application Housy}

\subsection{Interface Utilisateur}

\begin{enumerate}
    \item \textbf{Page d'accueil et dashboard principal}
    \begin{itemize}
        \item URL : \texttt{http://localhost:3000/dashboard}
        \item Description : Vue d'ensemble avec KPI, projets récents, notifications
        \item Éléments à capturer : Navigation, widgets, statistiques
    \end{itemize}
    
    \item \textbf{Module de gestion des projets}
    \begin{itemize}
        \item URL : \texttt{http://localhost:3000/projects}
        \item Description : Liste des projets, filtres, actions CRUD
        \item Éléments à capturer : Tableau, boutons d'action, pagination
    \end{itemize}
    
    \item \textbf{Formulaire de création de projet}
    \begin{itemize}
        \item URL : \texttt{http://localhost:3000/projects/new}
        \item Description : Wizard de création avec étapes multiples
        \item Éléments à capturer : Formulaire, validation, progression
    \end{itemize}
    
    \item \textbf{Interface d'estimation avec IA}
    \begin{itemize}
        \item URL : \texttt{http://localhost:3000/estimation}
        \item Description : Sélecteur de modèles IA, paramètres, résultats
        \item Éléments à capturer : AIModelSelector, calculs, justifications
    \end{itemize}
\end{enumerate}

\subsection{Fonctionnalités Avancées}

\begin{enumerate}
    \item \textbf{Sélecteur de modèles IA}
    \begin{itemize}
        \item Composant : \texttt{AIModelSelector.tsx}
        \item Description : Interface de choix entre 11 modèles disponibles
        \item Éléments à capturer : Badges, descriptions, recommandations
    \end{itemize}
    
    \item \textbf{Résultats d'estimation détaillés}
    \begin{itemize}
        \item Description : Breakdown des coûts par poste, justifications IA
        \item Éléments à capturer : Tableau détaillé, graphiques, sources
    \end{itemize}
    
    \item \textbf{Tableau de bord analytics}
    \begin{itemize}
        \item Description : KPI, métriques performance, tendances
        \item Éléments à capturer : Graphiques, indicateurs, filtres temporels
    \end{itemize}
    
    \item \textbf{Gestion des fournisseurs tunisiens}
    \begin{itemize}
        \item Description : Catalogue local, prix TND, disponibilité
        \item Éléments à capturer : Liste fournisseurs, fiches produits, tarifs
    \end{itemize}
\end{enumerate}

\section{Diagrammes Techniques Détaillés}

\subsection{Architecture Système}

\begin{enumerate}
    \item \textbf{Diagramme de déploiement Docker}
    \begin{itemize}
        \item Description : Containers, réseaux, volumes, dépendances
        \item Éléments : Frontend, Backend, Database, IA Services
        \item Format : Diagramme UML de déploiement
    \end{itemize}
    
    \item \textbf{Schéma de base de données}
    \begin{itemize}
        \item Description : Tables, relations, index, contraintes
        \item Éléments : ERD complet avec cardinalités
        \item Format : Diagramme entité-relation
    \end{itemize}
    
    \item \textbf{Pipeline CI/CD}
    \begin{itemize}
        \item Description : Étapes build, test, deploy automatisées
        \item Éléments : GitHub Actions, Docker Registry, déploiement
        \item Format : Diagramme de flux DevOps
    \end{itemize}
\end{enumerate}

\subsection{Intégration IA}

\begin{enumerate}
    \item \textbf{Architecture des modèles IA}
    \begin{itemize}
        \item Description : 11 modèles, service unifié, sélection automatique
        \item Éléments : Ollama, OpenAI, Perplexity, fallbacks
        \item Format : Diagramme d'architecture IA
    \end{itemize}
    
    \item \textbf{Flux d'enrichissement JSON}
    \begin{itemize}
        \item Description : Chargement données, enrichissement prompts
        \item Éléments : data-service, ai-service, estimation-service
        \item Format : Diagramme de flux de données
    \end{itemize}
    
    \item \textbf{Métriques et monitoring IA}
    \begin{itemize}
        \item Description : Tracking invisible, performance, qualité
        \item Éléments : Table tracking, KPI, dashboard développement
        \item Format : Dashboard métriques avec graphiques
    \end{itemize}
\end{enumerate}

\section{Recommandations pour la Finalisation}

\subsection{Priorités de Réalisation}

\begin{enumerate}
    \item \textbf{Priorité 1 - Indispensables}
    \begin{itemize}
        \item Logo entreprise ILOGsys
        \item Captures interface Housy principales
        \item Diagramme CRISP-DM adapté
        \item Architecture technique globale
    \end{itemize}
    
    \item \textbf{Priorité 2 - Importantes}
    \begin{itemize}
        \item Analyse comparative solutions existantes
        \item Chronogramme processus construction
        \item Captures fonctionnalités IA
        \item Schéma base de données
    \end{itemize}
    
    \item \textbf{Priorité 3 - Complémentaires}
    \begin{itemize}
        \item Infographies secteur construction
        \item Pipeline CI/CD
        \item Métriques IA détaillées
        \item Roadmap chapitres
    \end{itemize}
\end{enumerate}

\subsection{Outils Recommandés}

\begin{itemize}
    \item \textbf{Diagrammes BPMN :} Camunda Modeler, Lucidchart
    \item \textbf{Diagrammes UML :} PlantUML, StarUML, Draw.io
    \item \textbf{Captures d'écran :} Browser DevTools, Lightshot
    \item \textbf{Infographies :} Canva, Adobe Illustrator
    \item \textbf{Architecture :} Draw.io, Visio, Lucidchart
\end{itemize}

\begin{resultbox}
Ce guide fournit une liste exhaustive de toutes les illustrations à créer pour finaliser le rapport. Les diagrammes techniques principaux ont été générés en LaTeX/TikZ et sont prêts à l'utilisation. Les captures d'écran peuvent être réalisées directement depuis l'application Housy fonctionnelle.
\end{resultbox}
